\documentclass[aspectratio=169,10pt]{beamer}
\usetheme{Madrid}
\usecolortheme{seahorse}

\usepackage[utf8]{inputenc}
\usepackage{amsmath}
\usepackage{amssymb}
\usepackage{amsfonts}
\usepackage{graphicx}
\usepackage{listings}
\usepackage{xcolor}
\usepackage{booktabs}
\usepackage{array}
\usepackage{geometry}
\usepackage{fancyvrb}

\graphicspath{{figures/}}

% Define colors for code
\definecolor{codegreen}{rgb}{0.0, 0.5, 0.0}
\definecolor{codegray}{rgb}{0.5, 0.5, 0.5}
\definecolor{codepurple}{rgb}{0.58, 0, 0.82}
\definecolor{codeblue}{rgb}{0, 0, 0.8}

% Configure listings for Python code
\lstset{
    language=Python,
    basicstyle=\ttfamily\small,
    keywordstyle=\color{codeblue}\bfseries,
    commentstyle=\color{codegreen}\itshape,
    stringstyle=\color{codepurple},
    breaklines=true,
    numbers=left,
    numberstyle=\tiny\color{codegray},
    frame=single,
    rulecolor=\color{black},
    xleftmargin=2em,
    framexleftmargin=1.5em,
    showstringspaces=false
}

\setbeamertemplate{footline}[frame number]
\setbeamertemplate{navigation symbols}{}

\title[Chapter 5a: Fixed Point Theorems]{Chapter 5a: Fixed Point Theorems}
\subtitle{Metric Spaces, Contraction Mappings \& Applications}
\author{Based on Pathak's ``An Introduction to Nonlinear Analysis and Fixed Point Theory''}
\date{\today}
\institute{Nonlinear Analysis \& Fixed Point Theory}

\begin{document}

\begin{frame}
    \titlepage
\end{frame}

\begin{frame}{Outline}
    \tableofcontents
\end{frame}

\section{1. Introduction to Fixed Point Theory}

\begin{frame}{What is a Fixed Point?}
    \begin{definition}
        Let $T: X \to X$ be a mapping on a set $X$. A point $x^* \in X$ is called a \textbf{fixed point} of $T$ if
        \[
            T(x^*) = x^*.
        \]
    \end{definition}

    \vspace{0.5cm}
    \textbf{Key Question:} Under what conditions does a mapping $T$ have a fixed point?

    \vspace{0.5cm}
    \begin{example}
        Let $T: \mathbb{R} \to \mathbb{R}$ be defined by $T(x) = x + a$ for some fixed $a \neq 0$. \\
        Then $T$ has \textbf{no fixed point} since $T(x) = x$ implies $x + a = x$, a contradiction.
    \end{example}

    \vspace{0.3cm}
    \begin{example}
        Let $T: \mathbb{R} \to \mathbb{R}$ be defined by $T(x) = x^2$. \\
        Then $x = 0$ and $x = 1$ are fixed points.
    \end{example}
\end{frame}

\begin{frame}{Fixed Point Spaces}
    \begin{definition}
        A space $X$ is called a \textbf{fixed point space} if every continuous self-mapping of $X$ has a fixed point.
    \end{definition}

    \vspace{0.5cm}
    \textbf{Example:} $\mathbb{R}$ is \textit{not} a fixed point space (see $T(x) = x + 1$).

    \vspace{0.3cm}
    \begin{observation}
        Being a fixed point space is \textbf{topologically invariant}. If $X$ is a fixed point space and $Y$ is homeomorphic to $X$, then $Y$ is also a fixed point space.
    \end{observation}

    \vspace{0.5cm}
    \textbf{Key Insight:} Not every space has the fixed point property. Fixed point theory seeks to identify:
    \begin{itemize}
        \item Classes of spaces with the fixed point property
        \item Conditions on mappings ensuring fixed points exist
        \item Methods to compute or approximate fixed points
    \end{itemize}
\end{frame}

\begin{frame}[allowframebreaks]{Historical Perspective}
    Fixed point theory is divided into three major areas:

    \vspace{0.4cm}
    \textbf{1. Topological Fixed Point Theory}
    \begin{itemize}
        \item Involves topologically simple subsets of $\mathbb{R}^n$
        \item Mappings must be continuous
        \item Example: \textbf{Brouwer's Fixed Point Theorem} (1912)
            \begin{itemize}
                \item If $X$ is the unit ball in $\mathbb{R}^n$ and $T: X \to X$ is continuous, then $T$ has a fixed point
            \end{itemize}
    \end{itemize}

    \vspace{0.4cm}
    \textbf{2. Metric Fixed Point Theory}
    \begin{itemize}
        \item Operates in metric spaces with distance structure
        \item Focuses on contractive and nonexpansive mappings
        \item Example: \textbf{Banach Contraction Principle} (1922)
            \begin{itemize}
                \item Every complete metric space is a fixed point space for contraction mappings
            \end{itemize}
    \end{itemize}

    \vspace{0.4cm}
    \textbf{3. Lattice-Theoretic Fixed Point Theory}
    \begin{itemize}
        \item Uses order structure on spaces
        \item Works with partially ordered sets
        \item Applications to systems of inequalities
    \end{itemize}

    \framebreak

    \textbf{Historical Development:}
    \begin{itemize}
        \item \textbf{1912}: Brouwer proves the fixed point theorem for the unit ball in $\mathbb{R}^n$
        \item \textbf{1920}: Banach develops the contraction mapping principle
        \item \textbf{1930s-1960s}: Extensions to nonexpansive mappings, multifunctions
        \item \textbf{Modern Era}: Applications to optimization, differential equations, control theory
    \end{itemize}
\end{frame}

\section{2. Metric Spaces and Basic Concepts}

\begin{frame}{Metric Spaces}
    \begin{definition}
        A \textbf{metric space} $(X, d)$ consists of a set $X$ and a function $d: X \times X \to [0, \infty)$ such that for all $x, y, z \in X$:
        \begin{enumerate}
            \item \textbf{Non-negativity}: $d(x,y) \geq 0$ and $d(x,y) = 0 \Leftrightarrow x = y$
            \item \textbf{Symmetry}: $d(x,y) = d(y,x)$
            \item \textbf{Triangle Inequality}: $d(x,z) \leq d(x,y) + d(y,z)$
        \end{enumerate}
        The function $d$ is called a \textbf{metric} or \textbf{distance function}.
    \end{definition}

    \vspace{0.4cm}
    \textbf{Examples of Metric Spaces:}
    \begin{itemize}
        \item Euclidean space $(\mathbb{R}^n, \|\cdot\|_2)$ with $d(x,y) = \|x-y\|_2$
        \item Discrete metric: $d(x,y) = \begin{cases} 1 & \text{if } x \neq y \\ 0 & \text{if } x = y \end{cases}$
        \item Sequence spaces $\ell^p$ and $c_0$ with appropriate norms
    \end{itemize}

    \vspace{0.3cm}
    \begin{center}
        \includegraphics[width=0.7\textwidth]{fig_metric_space.pdf}
    \end{center}
\end{frame}

\begin{frame}[fragile]{Complete Metric Spaces}
    \begin{definition}
        A sequence $(x_n)$ in a metric space $(X, d)$ is \textbf{Cauchy} if for every $\epsilon > 0$, there exists $N$ such that
        \[
            n, m > N \Rightarrow d(x_n, x_m) < \epsilon.
        \]
    \end{definition}

    \vspace{0.3cm}
    \begin{definition}
        A metric space $(X, d)$ is \textbf{complete} if every Cauchy sequence in $X$ converges to some point in $X$.
    \end{definition}

    \vspace{0.4cm}
    \textbf{Examples:}
    \begin{itemize}
        \item $(\mathbb{R}^n, d_2)$ is complete (fundamental property of real numbers)
        \item The space $C[a,b]$ of continuous functions on $[a,b]$ with metric $d(f,g) = \max_{x \in [a,b]} |f(x) - g(x)|$ is complete
        \item The open interval $(0,1)$ with Euclidean metric is \textbf{not} complete
    \end{itemize}

    \vspace{0.3cm}
    \textbf{Key Fact:} Banach's contraction principle requires the underlying space to be complete!
\end{frame}

\begin{frame}{Normed Spaces and Banach Spaces}
    \begin{definition}
        A \textbf{norm} on a vector space $X$ is a function $\|\cdot\|: X \to [0, \infty)$ such that:
        \begin{enumerate}
            \item $\|x\| = 0 \Leftrightarrow x = 0$
            \item $\|\alpha x\| = |\alpha| \|x\|$ for all $\alpha \in \mathbb{R}$ (homogeneity)
            \item $\|x + y\| \leq \|x\| + \|y\|$ (triangle inequality)
        \end{enumerate}
        A \textbf{normed space} $(X, \|\cdot\|)$ is a vector space with a norm.
    \end{definition}

    \vspace{0.3cm}
    Every norm induces a metric: $d(x,y) = \|x - y\|$

    \vspace{0.4cm}
    \begin{definition}
        A \textbf{Banach space} is a complete normed vector space.
    \end{definition}

    \vspace{0.3cm}
    \textbf{Examples of Banach Spaces:}
    \begin{itemize}
        \item $(\mathbb{R}^n, \|\cdot\|_p)$ for any $p \in [1, \infty]$
        \item $(C[a,b], \|\cdot\|_\infty)$ where $\|f\|_\infty = \max_{x \in [a,b]} |f(x)|$
        \item $(\ell^p, \|\cdot\|_p)$ for $1 \leq p \leq \infty$
    \end{itemize}
\end{frame}

\section{3. Contractive Mappings}

\begin{frame}{Contractive Mappings}
    \begin{definition}
        Let $(X, d)$ be a metric space. A mapping $T: X \to X$ is called a \textbf{contraction} (or \textbf{contraction mapping}) if there exists $0 \leq \alpha < 1$ such that
        \[
            d(T(x), T(y)) \leq \alpha \, d(x, y) \quad \text{for all } x, y \in X.
        \]
        The constant $\alpha$ is called the \textbf{contraction coefficient}.
    \end{definition}

    \vspace{0.4cm}
    \textbf{Geometric Interpretation:} A contraction ``shrinks'' distances by at least a factor of $\alpha$. The closer $\alpha$ is to 0, the stronger the contraction.

    \vspace{0.4cm}
    \textbf{Observation:} Every contraction is uniformly continuous.

    \vspace{0.3cm}
    \begin{proof}[Proof Sketch]
        For $\epsilon > 0$, let $\delta = \epsilon/\alpha$. Then $d(x,y) < \delta$ implies
        \[
            d(T(x), T(y)) \leq \alpha \, d(x,y) < \alpha \delta = \epsilon.
        \]
    \end{proof}

    \vspace{0.4cm}
    \begin{center}
        \includegraphics[width=0.7\textwidth]{fig_contractive_mapping.pdf}
    \end{center}
\end{frame}

\begin{frame}{The Banach Contraction Principle}
    \begin{theorem}[Banach Contraction Principle]
        Let $(X, d)$ be a complete metric space and let $T: X \to X$ be a contraction with coefficient $0 \leq \alpha < 1$. Then:
        \begin{enumerate}
            \item $T$ has a \textbf{unique fixed point} $x^* \in X$
            \item For any starting point $x_0 \in X$, the sequence $(x_n)$ defined by $x_{n+1} = T(x_n)$ converges to $x^*$
            \item The rate of convergence is at least geometric: $d(x_n, x^*) \leq \alpha^n d(x_0, x^*)$
        \end{enumerate}
    \end{theorem}

    \vspace{0.4cm}
    \textbf{Significance:}
    \begin{itemize}
        \item Guarantees existence and uniqueness of the fixed point
        \item Provides an explicit iteration algorithm
        \item Gives bounds on the error and convergence rate
        \item One of the most practical fixed point results
    \end{itemize}

    \vspace{0.3cm}
    \begin{center}
        \includegraphics[width=0.65\textwidth]{fig_banach_theorem.pdf}
    \end{center}
\end{frame}

\begin{frame}[fragile]{Proof of Banach Contraction Principle}
    \textbf{Proof (outline):}

    \vspace{0.3cm}
    \textbf{Step 1: Cauchy sequence} \\
    For $x_n = T^n(x_0)$, we have
    \begin{align*}
        d(x_{n+1}, x_n) &= d(T(x_n), T(x_{n-1})) \leq \alpha \, d(x_n, x_{n-1}) \\
        &\leq \alpha^n \, d(x_1, x_0)
    \end{align*}
    For $m > n$:
    \[
        d(x_m, x_n) \leq (d(x_{n+1}, x_n) + d(x_{n+2}, x_{n+1}) + \cdots) \leq \frac{\alpha^n}{1-\alpha} d(x_1, x_0)
    \]
    So $(x_n)$ is Cauchy.

    \vspace{0.3cm}
    \textbf{Step 2: Convergence} \\
    Since $X$ is complete, $x_n \to x^*$ for some $x^* \in X$.

    \vspace{0.3cm}
    \textbf{Step 3: Fixed point} \\
    By continuity of $T$: $T(x^*) = T(\lim x_n) = \lim T(x_n) = \lim x_{n+1} = x^*$.

    \vspace{0.3cm}
    \textbf{Step 4: Uniqueness} \\
    If $T(y) = y$, then $d(x^*, y) = d(T(x^*), T(y)) \leq \alpha \, d(x^*, y)$, so $(1-\alpha) d(x^*, y) \leq 0$, thus $x^* = y$.
\end{frame}

\begin{frame}[fragile]{Examples of Contractions}
    \textbf{Example 1:} $T: [0,1] \to [0,1]$ defined by $T(x) = 0.5x + 1$ \\
    The contraction coefficient is $\alpha = 0.5 < 1$. Fixed point: $x^* = 2$.

    \vspace{0.3cm}
    \begin{lstlisting}[language=Python]
# Compute fixed point iteratively
x = 5.0
for i in range(20):
    x = 0.5 * x + 1
    print(f"Iteration {i+1}: x = {x:.6f}")
# Output converges to x* = 2
    \end{lstlisting}

    \vspace{0.3cm}
    \textbf{Example 2:} Integral equation
    \[
        x(t) = f(t) + \int_0^t K(t,s) x(s) \, ds
    \]
    Under suitable conditions on $K$ and $f$, the operator
    \[
        T(x)(t) = f(t) + \int_0^t K(t,s) x(s) \, ds
    \]
    is a contraction on $C[0, T]$, guaranteeing a unique solution.

    \vspace{0.3cm}
    \textbf{Example 3:} Differential equations \\
    The initial value problem $\frac{dx}{dt} = f(t,x)$, $x(0) = x_0$ can be converted to the integral equation
    \[
        x(t) = x_0 + \int_0^t f(s, x(s)) \, ds
    \]
    If $f$ is Lipschitz continuous, this generates a contraction, ensuring existence and uniqueness of solutions.
\end{frame}

\section{4. Nonexpansive Mappings}

\begin{frame}{Nonexpansive Mappings}
    \begin{definition}
        A mapping $T: X \to X$ on a metric space $(X, d)$ is called \textbf{nonexpansive} if
        \[
            d(T(x), T(y)) \leq d(x, y) \quad \text{for all } x, y \in X.
        \]
    \end{definition}

    \vspace{0.4cm}
    \textbf{Comparison:}
    \begin{itemize}
        \item \textbf{Contraction}: $d(T(x), T(y)) \leq \alpha \, d(x,y)$ with $\alpha < 1$
        \item \textbf{Nonexpansive}: $d(T(x), T(y)) \leq d(x,y)$ (does \textit{not} shrink distances)
    \end{itemize}

    \vspace{0.4cm}
    \textbf{Key Difference:} Nonexpansive mappings do not necessarily have fixed points! For example:
    \[
        T: [0,1] \to [0,1], \quad T(x) = x + 1 \pmod{1}
    \]
    is nonexpansive but has no fixed point.

    \vspace{0.3cm}
    \textbf{Fixed Points for Nonexpansive Mappings:}
    \begin{itemize}
        \item \textbf{Theorem (Kirk, 1965)}: Every nonexpansive self-mapping of a weakly compact convex subset of a Banach space with the fixed point property for nonexpansive mappings has a fixed point
        \item The existence requires additional assumptions (compactness, convexity, reflexivity)
    \end{itemize}
\end{frame}

\begin{frame}{Examples of Nonexpansive Mappings}
    \textbf{Example 1: Projection Operator} \\
    Let $K$ be a closed convex subset of a Hilbert space. The projection onto $K$,
    \[
        P_K(x) = \arg\min_{y \in K} \|x - y\|,
    \]
    is nonexpansive. Fixed points are exactly the points in $K$.

    \vspace{0.4cm}
    \textbf{Example 2: Averaging Operator} \\
    In $\mathbb{R}^n$, define $T(x) = \frac{1}{n}(x + x_1) \cdot \mathbf{1}$. This is nonexpansive.

    \vspace{0.4cm}
    \textbf{Example 3: Metric Contraction} \\
    Let $T: X \to X$ satisfy $\|T(x) - T(y)\| \leq L \|x - y\|$ with $L \leq 1$. This is nonexpansive.

    \vspace{0.4cm}
    \textbf{Example 4: Composition} \\
    If $T_1, T_2$ are nonexpansive, then their composition $T_1 \circ T_2$ is also nonexpansive.
\end{frame}

\section{5. Applications and Examples}

\begin{frame}[fragile]{Solving Linear Systems}
    Consider the linear system $Ax = b$ where $A \in \mathbb{R}^{n \times n}$ is invertible.

    \vspace{0.3cm}
    \textbf{Iterative approach:} Define $T(x) = x + A^{-1}(b - Ax) = x - A^{-1}(Ax - b)$

    \vspace{0.3cm}
    Then $T(x) = x \Leftrightarrow A^{-1}(Ax - b) = 0 \Leftrightarrow x$ solves $Ax = b$.

    \vspace{0.3cm}
    For convergence, we need $T$ to be a contraction. This happens when the spectral radius of $(I - A^{-1}A) = 0$... Actually, let's use the Richardson iteration:
    \[
        T(x) = (I - \alpha A) x + \alpha b
    \]
    with $\alpha$ chosen so that $\|(I - \alpha A)\| < 1$.

    \vspace{0.3cm}
    \begin{lstlisting}[language=Python]
import numpy as np
A = np.array([[2, 1], [1, 3]], dtype=float)
b = np.array([5, 6], dtype=float)
alpha = 0.1  # Step size

x = np.array([0, 0], dtype=float)
for i in range(50):
    x = (np.eye(2) - alpha*A) @ x + alpha * b
    print(f"Iteration {i+1}: x = {x}")
    \end{lstlisting}
\end{frame}

\begin{frame}{Applications to Differential Equations}
    \textbf{Picard's Iteration Theorem} \\
    Consider the initial value problem:
    \[
        \frac{dx}{dt} = f(t, x), \quad x(0) = x_0
    \]
    where $f$ is continuous and satisfies a Lipschitz condition in $x$:
    \[
        |f(t, x) - f(t, y)| \leq L |x - y|.
    \]

    \vspace{0.3cm}
    \textbf{Equivalent integral equation:}
    \[
        x(t) = x_0 + \int_0^t f(s, x(s)) \, ds
    \]

    \vspace{0.3cm}
    \textbf{Picard's iteration:}
    \[
        x_{n+1}(t) = x_0 + \int_0^t f(s, x_n(s)) \, ds
    \]

    \vspace{0.3cm}
    On the space $C[0, T]$ with the metric $d(x, y) = \sup_{t \in [0,T]} |x(t) - y(t)|$, the operator
    \[
        T(x)(t) = x_0 + \int_0^t f(s, x(s)) \, ds
    \]
    becomes a contraction for sufficiently small $T > 0$ (depending on $L$). Thus, there exists a unique solution on $[0, T]$.
\end{frame}

\begin{frame}[fragile]{Numerical Example: Newton-Kantorovich Method}
    \textbf{Newton-Kantorovich Method for $f(x) = 0$:}
    \[
        T(x) = x - \frac{f(x)}{f'(x)}
    \]

    \vspace{0.3cm}
    If $x^*$ is a simple root ($f(x^*) = 0$ and $f'(x^*) \neq 0$), then $T$ is a contraction near $x^*$ with quadratic convergence!

    \vspace{0.3cm}
    \begin{lstlisting}[language=Python]
# Newton's method for f(x) = x^3 - 2
def f(x):
    return x**3 - 2

def f_prime(x):
    return 3 * x**2

x = 1.0  # Initial guess
for i in range(10):
    x_new = x - f(x) / f_prime(x)
    error = abs(x_new - x)
    print(f"Iter {i+1}: x = {x_new:.10f}, error = {error:.2e}")
    x = x_new
# Converges rapidly to x^* = 2^(1/3)
    \end{lstlisting}
\end{frame}

\section{6. Convergence Analysis}

\begin{frame}{Convergence Rates}
    \begin{definition}
        Let $(x_n)$ be a sequence converging to $x^*$. Define $e_n = d(x_n, x^*)$.
        \begin{enumerate}
            \item \textbf{Linear convergence}: There exists $0 < \rho < 1$ and $C > 0$ such that $e_{n+1} \leq C \rho^n e_0$
            \item \textbf{Superlinear convergence}: $\lim_{n \to \infty} \frac{e_{n+1}}{e_n} = 0$
            \item \textbf{Quadratic convergence}: There exists $C > 0$ such that $e_{n+1} \leq C e_n^2$
        \end{enumerate}
    \end{definition}

    \vspace{0.4cm}
    \textbf{Banach Contraction Analysis:}
    For contraction with coefficient $\alpha$:
    \[
        d(x_n, x^*) \leq \alpha^n d(x_0, x^*)
    \]
    This is linear convergence with rate $\rho = \alpha$.

    \vspace{0.3cm}
    \textbf{Number of iterations for accuracy $\epsilon$:}
    \[
        \alpha^n d(x_0, x^*) < \epsilon \Rightarrow n > \frac{\log(\epsilon / d(x_0, x^*))}{\log \alpha}
    \]

    \vspace{0.3cm}
    \begin{center}
        \includegraphics[width=0.7\textwidth]{fig_convergence_behavior.pdf}
    \end{center}
\end{frame}

\begin{frame}{Cobweb Diagrams}
    \textbf{Visualization of Fixed Point Iteration:}

    A cobweb diagram shows graphically how iterations $x_{n+1} = T(x_n)$ converge to the fixed point.

    \vspace{0.3cm}
    \textbf{Construction:}
    \begin{enumerate}
        \item Plot $y = T(x)$ and $y = x$
        \item Start at $(x_0, 0)$ on the x-axis
        \item Move vertically to $(x_0, T(x_0))$
        \item Move horizontally to $(T(x_0), T(x_0))$ on the line $y = x$
        \item Repeat, creating a cobweb pattern
    \end{enumerate}

    \vspace{0.3cm}
    \textbf{Convergence behavior:}
    \begin{itemize}
        \item If $|T'(x^*)| < 1$ (contraction): Cobweb spirals toward the fixed point
        \item If $|T'(x^*)| > 1$ (expansion): Cobweb diverges from the fixed point
        \item If $|T'(x^*)| = 1$: Behavior depends on higher derivatives
    \end{itemize}

    \vspace{0.3cm}
    \begin{center}
        \includegraphics[width=0.7\textwidth]{fig_cobweb_diagram.pdf}
    \end{center}
\end{frame}

\section{7. Extended Results}

\begin{frame}[allowframebreaks]{Generalizations of Banach's Theorem}
    \textbf{Kannan Fixed Point Theorem}
    \begin{itemize}
        \item For complete metric spaces, Kannan contractions also guarantee fixed points
        \item A mapping $T$ is Kannan if $d(T(x), T(y)) \leq \alpha(d(x, T(x)) + d(y, T(y)))$ with $\alpha < 1/2$
        \item More general than Banach contractions
    \end{itemize}

    \vspace{0.3cm}
    \textbf{Chatterjea Fixed Point Theorem}
    \begin{itemize}
        \item Even more general class of contractions
        \item Applies to complete metric spaces
        \item Useful for mappings that satisfy $d(T(x), T(y)) \leq \alpha(d(x, T(y)) + d(y, T(x)))$
    \end{itemize}

    \vspace{0.3cm}
    \textbf{Browder Fixed Point Theorem}
    \begin{itemize}
        \item Every nonexpansive self-mapping of a bounded closed convex subset of a uniformly convex Banach space has a fixed point
        \item Applies to nonexpansive mappings (not just contractions)
        \item Requires reflexivity and uniform convexity
    \end{itemize}

    \framebreak

    \textbf{Kirk Fixed Point Theorem}
    \begin{itemize}
        \item Extends results to hyperbolic metric spaces
        \item Every nonexpansive mapping of a weakly compact subset has a fixed point
        \item Bridges topological and metric fixed point theory
    \end{itemize}

    \vspace{0.4cm}
    \begin{center}
        \includegraphics[width=0.8\textwidth]{fig_theorems_comparison.pdf}
    \end{center}
\end{frame}

\begin{frame}{Strong Contractions (Pathak-Shahzad, 2008)}
    \begin{definition}
        A mapping $T: X \to X$ on a metric space is called a \textbf{strong contraction} if
        \begin{enumerate}
            \item $T$ is a contraction with some $\alpha < 1$
            \item $\delta(T^n(X)) \to 0$ as $n \to \infty$ (diameter of iterates shrinks)
        \end{enumerate}
    \end{definition}

    \vspace{0.3cm}
    \textbf{Theorem (Pathak-Shahzad):} Every strong topological contraction mapping on a topologically complete metric space endowed with metric topology has a unique fixed point $w$ such that $\{w\} = \bigcap_{n=0}^\infty T^n(X)$.

    \vspace{0.3cm}
    \textbf{Key Results:}
    \begin{itemize}
        \item Existence and uniqueness of the fixed point
        \item Provides a characterization: the fixed point is the nested intersection of iterates
        \item Extends classical Banach principle to more general settings
        \item Works in topologically complete spaces (weaker than metric completeness)
    \end{itemize}
\end{frame}

\begin{frame}{Strongly $p$-Contractive Mappings}
    \begin{definition}
        Let $(X, d)$ be a metric space and $T: X \to X$. We say $T$ is \textbf{strongly $p$-contractive} if $T^p$ (the $p$-fold composition of $T$) is a contraction for some $p \in \mathbb{N}$.
    \end{definition}

    \vspace{0.3cm}
    \textbf{Theorem (Pathak-Shahzad):} Let $(X, d)$ be a metric space and let $T: X \to X$ be a strongly $p$-contraction mapping. If $T$ is orbitally complete, then $T$ has a unique fixed point.

    \vspace{0.3cm}
    \textbf{Definition (Orbital Completeness):} A metric space is \textbf{orbitally complete} with respect to a mapping $T$ if every Cauchy sequence consisting of iterates $\{T^n(x)\}$ converges in $X$.

    \vspace{0.3cm}
    \textbf{Implications:}
    \begin{itemize}
        \item Weaker than metric completeness
        \item Allows for fixed points even when the entire space is not complete
        \item Useful for establishing fixed points of power mappings
    \end{itemize}
\end{frame}

\section{8. Summary and Key Takeaways}

\begin{frame}{Key Theorems Summary}
    \begin{table}
        \centering
        \small
        \begin{tabular}{|l|c|c|c|}
            \hline
            \textbf{Theorem} & \textbf{Space} & \textbf{Mapping} & \textbf{Fixed Point?} \\
            \hline
            Brouwer & Convex, compact & Continuous & Yes \\
            \hline
            Banach & Complete metric & Contraction & Yes, unique \\
            \hline
            Browder & Unif. convex Banach & Nonexpansive & Yes \\
            \hline
            Schauder & Banach & Compact & Yes \\
            \hline
            Kirk & Hyperbolic space & Nonexpansive & Yes \\
            \hline
            Kannan & Complete metric & Kannan contract. & Yes, unique \\
            \hline
        \end{tabular}
    \end{table}

    \vspace{0.4cm}
    \textbf{Unifying Theme:} Different fixed point theorems trade off:
    \begin{itemize}
        \item Space structure (convexity, completeness, compactness)
        \item Mapping properties (continuity, contraction, nonexpansiveness)
    \end{itemize}
\end{frame}

\begin{frame}{Practical Applications}
    \textbf{1. Optimization and Numerical Analysis}
    \begin{itemize}
        \item Iterative algorithms (gradient descent, Newton's method)
        \item Convergence analysis and rates
        \item Error bounds and stopping criteria
    \end{itemize}

    \vspace{0.3cm}
    \textbf{2. Differential Equations}
    \begin{itemize}
        \item Existence and uniqueness of solutions to ODEs and PDEs
        \item Integral equations (Volterra, Fredholm)
        \item Boundary value problems
    \end{itemize}

    \vspace{0.3cm}
    \textbf{3. Game Theory and Economics}
    \begin{itemize}
        \item Nash equilibrium as fixed points
        \item Repeated games and stationary equilibria
        \item Convergence to equilibrium
    \end{itemize}

    \vspace{0.3cm}
    \textbf{4. Image Processing and Computer Vision}
    \begin{itemize}
        \item Fixed point iterations in image denoising
        \item Proximal algorithms for convex optimization
        \item Convergence guarantees for algorithms
    \end{itemize}
\end{frame}

\begin{frame}[fragile]{Computational Considerations}
    \textbf{When choosing an iterative method:}

    \vspace{0.3cm}
    \textbf{Algorithm Selection:}
    \begin{lstlisting}[language=Python]
# Check contraction coefficient
alpha = estimate_contraction_coeff(T, domain)
if alpha < 0.1:
    # Use simple fixed point iteration
    x = fixed_point_iteration(T, x0)
elif alpha < 0.5:
    # Good contraction, maybe accelerate
    x = aitken_acceleration(T, x0)
elif alpha < 1.0:
    # Weak contraction, consider alternatives
    x = anderson_acceleration(T, x0)
else:
    # Not a contraction, use Newton or quasi-Newton
    x = newton_method(f, x0)
    \end{lstlisting}

    \vspace{0.3cm}
    \textbf{Convergence Monitoring:}
    \begin{itemize}
        \item Track $\|x_{n+1} - x_n\|$ for convergence detection
        \item Use error bounds from contraction coefficient
        \item Implement adaptive step sizes
    \end{itemize}
\end{frame}

\begin{frame}{Further Reading and Resources}
    \textbf{Foundational Works:}
    \begin{itemize}
        \item Banach, S. (1922). "Sur les opérations dans les ensembles abstraits et leur application aux équations intégrales"
        \item Brouwer, L.E.J. (1912). "Über Abbildung von Mannigfaltigkeiten"
        \item Kirk, W.A. (1965). "A fixed point theorem for mappings which do not increase distances"
    \end{itemize}

    \vspace{0.3cm}
    \textbf{Modern References:}
    \begin{itemize}
        \item Pathak, H.K. (2018). "An Introduction to Nonlinear Analysis and Fixed Point Theory"
        \item Granas, A. \& Dugundji, J. (2003). "Fixed Point Theory"
        \item Smart, D.R. (1974). "Fixed Point Theorems"
    \end{itemize}

    \vspace{0.3cm}
    \textbf{Application Areas:}
    \begin{itemize}
        \item Numerical linear algebra and optimization
        \item Functional analysis and operator theory
        \item PDEs and integral equations
        \item Machine learning and optimization
    \end{itemize}
\end{frame}

\begin{frame}{Conclusion}
    \begin{center}
        {\Large \textbf{Fixed Point Theorems}} \\
        \vspace{0.5cm}
        A foundational tool bridging abstract mathematics and practical computation
    \end{center}

    \vspace{0.5cm}
    \textbf{Key Insights:}
    \begin{itemize}
        \item Not every mapping has a fixed point, but contractions always do
        \item Completeness and contraction together guarantee existence and uniqueness
        \item Iteration provides both a constructive proof and an algorithm
        \item Generalizations extend to broader classes of spaces and mappings
    \end{itemize}

    \vspace{0.5cm}
    \textbf{Impact:}
    \begin{itemize}
        \item Fundamental in nonlinear analysis
        \item Essential for numerical methods
        \item Critical for understanding equilibrium and convergence
        \item Continues to drive research in functional analysis
    \end{itemize}

    \vspace{0.8cm}
    \begin{center}
        {\large Thank you!}
    \end{center}
\end{frame}

\end{document}
